%%%%%%%%%%%%%%%%%%%%%%%%%%%%%%%%%%%%%%%%%%%%%%%%%
%%%%%%%%%%%% Week 1: Introduction %%%%%%%%%%%%%%%
%%%%%%%%%%%%%%%%%%%%%%%%%%%%%%%%%%%%%%%%%%%%%%%%%

\chapter{Introduction}\label{chap:intro}

% STATUS: Week 1 (completed)

\section{Motivation}\label{sec:motivation}

Detectors trained on GAN (Generative Adversarial Network) images have been shown to produce near-random accuracy when tested on diffusion model outputs. With the public release of generative models -- including GANs, latent diffusion models, and autoregressive transformers -- the divide between image generation capabilities and detection has grown exponentially, with real-world impacts on forensics, journalism, and law.

Existing detectors suffer from three key limitations:
\begin{enumerate}
    \item \textbf{Lack of calibrated uncertainty}: Most detectors output a binary real/fake label without any confidence estimate. When wrong, they are confidently wrong.
    \item \textbf{Fragility under compression}: Social-media platforms (WhatsApp, Instagram) recompress uploaded images, degrading forensic artifacts that detectors rely upon. Prior work reports 3--8\% AUC loss under recompression.
    \item \textbf{Poor cross-generator generalisation}: A detector trained on one generator family (e.g., ProGAN) fails when evaluated on unseen architectures (e.g., Stable Diffusion, DALL-E).
\end{enumerate}

To address these limitations, we propose \textbf{ImageTrust}, an image forensic framework that combines multi-backbone feature fusion with calibrated uncertainty quantification.

\section{Objectives}\label{sec:objectives}

The objectives of this thesis are:
\begin{enumerate}
    \item \textbf{Calibrated uncertainty quantification}: Combine temperature scaling (to minimise calibration error) with conformal prediction (to provide coverage guarantees with explicit abstention). To our knowledge, this combination has not been studied in existing AI-generated image detectors.
    \item \textbf{Compression-robust multi-backbone fusion}: Fuse ResNet-50, EfficientNet-B0, and ViT-B/16 embeddings into a 4,096-dimensional vector, classified by XGBoost and MLP meta-classifiers, achieving less than 0.3\% AUC degradation under social-media compression.
    \item \textbf{End-to-end forensic system}: Integrate forensic modules (copy-move detection, C2PA provenance, Grad-CAM explainability) into a deployable desktop and web application.
    \item \textbf{Rigorous evaluation}: Benchmark against six baselines (four internal, two external) on 604,589 images, four compression variants, 24 unseen generators, and a component ablation study.
\end{enumerate}

\section{Thesis Structure}\label{sec:structure}

This thesis is organised as follows:
\begin{itemize}
    \item \textbf{Chapter~\ref{chap:data}} presents the dataset construction, data splitting strategy, training configuration, baselines definition, and main comparison results.
    \item \textbf{Chapter~\ref{chap:related}} reviews the state of the art in AI-generated image detection, calibration, and conformal prediction.
    \item \textbf{Chapter~\ref{chap:methodology}} describes the methodology: heterogeneous backbone fusion, meta-classifier architectures, temperature scaling calibration, and conformal prediction.
    \item \textbf{Chapter~\ref{chap:results}} presents detailed experimental results: degradation robustness, cross-generator generalisation, ablation study, calibration analysis, and efficiency benchmarks.
    \item \textbf{Chapter~\ref{chap:system}} describes the end-to-end forensic system: desktop application (PySide6), web application (Next.js/FastAPI), and forensic modules.
    \item \textbf{Chapter~\ref{chap:conclusion}} concludes the thesis with a summary, limitations, and future work.
\end{itemize}
